\section{Introduction}\label{sec:introduction}

At a fixed temperature below the transition, how large are the
sample-to-sample fluctuations of the Sherrington--Kirkpatrick free
energy, and do they approach a law after centering and rescaling?
We prove a limiting law for every fixed inverse temperature $\beta>1$
in the zero-field Gaussian Ising model. The result includes a positive
variance prefactor at scale $n^{1/3}$ and a characterization of the
law by finite-system moments.

Let $(g_{ij})_{1\le i<j\le n}$ be independent standard Gaussian variables. For $\sigma\in\{-1,1\}^n$ and $\beta>1$, put
\begin{equation}\label{eq:model}
 H_n(\sigma)=\frac1{\sqrt n}\sum_{i<j}g_{ij}\sigma_i\sigma_j,
 \qquad Z_n(\beta)=\sum_{\sigma\in\{-1,1\}^n}e^{\beta H_n(\sigma)},
 \qquad F_n(\beta)=\log Z_n(\beta).
\end{equation}
All expectations concern the Gaussian disorder unless another probability law is explicitly specified. There is no external field. Write
\[
 s_n(\beta)^2=\Var F_n(\beta),
 \qquad X_n(\beta)=\frac{F_n(\beta)-\mathbb E F_n(\beta)}{s_n(\beta)}.
\]
The variance is positive for every $n\ge2$: with all couplings except $g_{12}$ set to zero, the free energy is $n\log2+\log\cosh(\beta g_{12}/\sqrt n)$, a nonconstant continuous function. Full support of the Gaussian law excludes almost-sure constancy. Standard Gaussian moment bounds give finite variance.

\subsection{The fluctuation problem and its history}

The model \eqref{eq:model} was introduced by Sherrington and
Kirkpatrick~\cite{SK1975}. Their symmetric replica calculation gave
an unphysical negative entropy at zero temperature. Parisi's
hierarchy~\cite{Parisi1979} replaced a single overlap parameter by
successive levels, encoded in the infinite-level limit by a function.
This development led to a variational formula for the thermodynamic
free energy.
Gaussian interpolation established the thermodynamic limit
\cite{GuerraToninelli2002}, the variational bound of
Guerra~\cite{Guerra2003}, and ultimately the matching formula proved by
Talagrand~\cite{Talagrand2006}. Panchenko's extension treats general
mixed $p$-spin models, including odd interactions~\cite{Panchenko2014}. These results identify the
leading order of the mean free energy. They leave the smaller,
sample-dependent fluctuations studied here.

The one-sixth fluctuation prediction has a separate physics history.
Kondor's replica expansion~\cite{Kondor1983} and the finite-size
prediction of Crisanti, Paladin, Sommers, and Vulpiani~\cite{CPSV1992}
preceded the low-temperature large-deviation calculations of Parisi
and Rizzo~\cite{ParisiRizzo2008,ParisiRizzo2009}. The predicted
$n^{-5/6}$ scale for the physical free-energy density becomes
$n^{1/6}$ for $\log Z_n$ at fixed $\beta$. To pass from rare
free energies to the width of typical fluctuations, Parisi and Rizzo
assume that the large-deviation tail matches smoothly onto the
central part of the distribution~\cite[Section~VI]{ParisiRizzo2009}.
They distinguish the two orders of the replica-number and system-size
limits. This matching argument is heuristic: it does not itself
provide a variance asymptotic or a limiting law.

Rigorous fluctuation results depend strongly on the regime.
Aizenman, Lebowitz, and Ruelle proved order-one Gaussian fluctuations
at high temperature~\cite{AizenmanLebowitzRuelle1987}.
Chatterjee obtained superconcentration and its relation to disorder
chaos~\cite{Chatterjee2009}, as well as a constant-scale
nonconcentration bound~\cite{Chatterjee2019LowerBounds}.
Chen and Lam's critical and near-critical bounds concern temperatures
approaching the transition~\cite{ChenLam2019Critical}.
For each fixed $\beta>1$, Aronow and Lopatto prove
\[
 c_\beta n^{4/15}\le \Var F_n(\beta)\le C_\beta n^{7/15}
\]
in \cite[Theorem~1.4]{AL2026}, and quantitative tail estimates in a
separate deviation window in \cite[Theorem~1.5]{AL2026}.
These bounds allow, but do not establish, variance of order $n^{1/3}$.
Recent critical-window limits~\cite[Theorem~2.1]{ChengLiuShaoXu2026CriticalWindow}
concern $\beta_n\to1$, while the all-temperature results of
Chen, Dey, and Panchenko require a nonzero field~\cite{ChenDeyPanchenko2017}.
Neither hypothesis covers the fixed zero-field model here.

Our companion \emph{The low-temperature Sherrington--Kirkpatrick
fluctuation scale}~\cite{PE}, abbreviated \emph{PE}, supplies both its
finite-tree comparison toolkit and the fixed-temperature exponent
\begin{equation}\label{eq:PE-exponent}
 s_n(\beta)=n^{1/6+o(1)}.
\end{equation}
The exponent alone does not identify a limiting distribution or prove
convergence of the variance divided by $n^{1/3}$. The additional work
here is a quantitative comparison of different sizes, with a summable
error. Appendix~\ref{sec:companion-interface} lists the precise PE
inputs and their hypotheses.

\subsection{The limiting law}

\begin{theorem}[Limiting law]\label{thm:main}
For every fixed real $\beta>1$, there is a nondegenerate Borel probability measure $\nu_\beta$ such that
\[
 X_n(\beta)\ \Longrightarrow\ \nu_\beta
 \qquad\text{as }n\to\infty\text{ through all integers}.
\]
The measure $\nu_\beta$ has mean zero, variance one, and a finite exponential moment in a neighborhood of the origin. It is the unique probability measure with the moments specified in Proposition~\ref{def:law} below.
\end{theorem}

The centering and variance are the exact finite-size quantities. The
formula below specifies the moments of $\nu_\beta$; no identification with a named
distribution is made. The quantitative statement
below proves that the variance has a positive limiting prefactor and
provides the rate needed for the moment characterization.

For the quantitative result, we use a Gaussian completion convenient for interpolation. Let $\zeta_n$ be a centered Gaussian of variance $\beta^2/2$, independent of the couplings, and define
\begin{equation}\label{eq:completed-free-energy}
 F_n^c=F_n(\beta)-n\log2+\zeta_n,
 \qquad Y_n=n^{-1/6}(F_n^c-\mathbb E F_n^c).
\end{equation}
Different copies below include independent couplings and independent completion variables. The subtraction of $n\log2$ changes the counting spin prior into the uniform probability prior. The completion changes the covariance into the homogeneous quadratic-overlap covariance used later.

\begin{theorem}[Moment convergence at the fluctuation scale]\label{thm:moments}
For each fixed $\beta>1$, there are real numbers $a_k=a_k(\beta)$, $k\ge1$, with the following property. For every integer $k\ge1$, there are constants $c_k>0$ and $C_k<\infty$ such that
\begin{equation}\label{eq:moment-rate}
 \big|\mathbb E Y_n^k-a_k\big|\le C_k n^{-c_k},\qquad n\ge2.
\end{equation}
There are also constants $c>0$ and $C<\infty$, depending on $\beta$, for which
\begin{equation}\label{eq:exponential-moments}
 \sup_{n\ge2}\mathbb E e^{c|Y_n|}\le C.
\end{equation}
The limits satisfy $a_1=0$ and $0<a_2<\infty$.
\end{theorem}

Theorem~\ref{thm:main} follows from Theorem~\ref{thm:moments} by moment determinacy and removal of the completion. Indeed,
\[
 n^{-1/3}\Var F_n(\beta)
 =\Var Y_n-\frac{\beta^2}{2n^{1/3}}\longrightarrow a_2>0,
\]
and the centered completion vanishes in $L^2$ after multiplication by $n^{-1/6}$. The exact finite-size standardization is then justified by Slutsky's theorem. We give the complete deduction, including nondegeneracy, in Section~\ref{sec:reduction}.
For the physical total free energy $-\beta^{-1}F_n$, the variance
constant is $a_2/\beta^2$ and its standardized limit is the reflection
of $\nu_\beta$. For its density $-(\beta n)^{-1}F_n$, the variance is
asymptotic to $(a_2/\beta^2)n^{-5/3}$.

\subsection{A finite-system characterization of the law}

The power rate in Theorem~\ref{thm:moments} allows each limiting
moment to be recovered from a convergent series of finite-size averages.
The following random-size experiment expresses this characterization
using independent copies of the finite SK model.

\begin{proposition}[Moment specification]\label{def:law}
Fix $\beta>1$ and let $C=\sum_{l=2}^{\infty}l^{-2}$. For an integer $k\ge1$, draw an integer $J\ge2$ with
\[
 \mathbb P(J=l)=\frac{l^{-2}}C.
\]
Conditional on $J$, draw $k+1$ independent copies $F_J^{c,0},\ldots,F_J^{c,k}$ of the completed free energy in \eqref{eq:completed-free-energy}. Set
\[
 t=\log J,
 \qquad P_k=\prod_{i=1}^k J^{-1/6}(F_J^{c,0}-F_J^{c,i}).
\]
Conditional averaging over the independent comparison copies gives
$\mathbb E[P_k\mid J=l]=\mathbb E Y_l^k$.
The constants in Theorem~\ref{thm:moments} satisfy
\begin{equation}\label{eq:law-moments}
 a_k=C\sum_{h=0}^{\infty}
 \mathbb E\!\left[P_k\left(\frac{t^h}{h!}
       -\mathbf1_{\{h\ge1\}}\frac{t^{h-1}}{(h-1)!}\right)\right],
\end{equation}
where the second term is absent for $h=0$. Each expectation is taken before the series is summed. The series of these scalar expectations is absolutely convergent. We have $a_2>0$, and the law in Theorem~\ref{thm:main} is uniquely determined by
\[
 \int z^k\,\nu_\beta(dz)=\frac{a_k}{a_2^{k/2}},\qquad k\ge1.
\]
\end{proposition}

Section~\ref{sec:reduction} proves the formula and its absolute convergence from the moment rate in Theorem~\ref{thm:moments}.

This characterization involves only a fixed prior on finite system sizes and finite Gaussian experiments. It does not require the choice of a subsequence or knowledge of an asymptotic variance. The order in \eqref{eq:law-moments} is essential: it is not permissible to interchange the unaveraged pointwise series with expectation. Section~\ref{sec:reduction} also realizes each $a_k$ as the expectation of a single integrable random variable obtained from an almost surely finite experiment.

\subsection{The size comparison and its main ingredients}

The new estimate compares centered log Laplace transforms at different
system sizes with a summable error. For a large dyadic size $N$, put
$x=N^{-1/6}$. The comparison gives, for every integer $N'\in[N,2N]$,
\[
 \left|\log\mathbb E e^{q x^\varepsilon Y_{N'}}
       -\log\mathbb E e^{q x^\varepsilon Y_N}\right|
 \le C_q x^{\varepsilon+c'},
\]
where $c'>0$, $q>0$ is fixed, and $\varepsilon>0$ is sufficiently
small and fixed. The saving $c'$ is independent of sufficiently small
$\varepsilon$. Estimates for positive Laplace tests, together with
preliminary exponential-tail bounds, recover comparisons of each fixed
moment. Their errors can then be summed over dyadic intervals. This is the reduction
proved in Section~\ref{sec:reduction}.

To obtain the transform estimate, we interpolate through hierarchical
Gaussian systems. At each level, a normalized transition kernel
specifies the law of the next Gaussian increments; paths can share an
initial history and then branch independently. Overlaps are normalized
scalar products of the terminal spin configurations. The hierarchy is
indexed by an averaging parameter, called mass; its matrix and field
clocks record the cumulative variances along this index. A one-site
scalar model built from the Parisi minimizer supplies reference overlaps
along these clocks.

The scalar input has its own development. Auffinger and Chen proved
structural properties of Parisi measures and uniqueness of the Parisi minimizer
\cite{AuffingerChenProperties2015,AuffingerChenUniqueness2015};
Jagannath and Tobasco gave variational support conditions
\cite{JagannathTobasco2017}. Zhou established interval support near the
critical temperature~\cite{Zhou2025}. We use Lopatto's full-support
result for every fixed $\beta>1$~\cite[Theorem~1.1 and Lemma~2.2]{Lopatto2026},
in the form developed in PE. The positive initial density needed here
is proved in that companion from the scalar identities and full
support.

The direct antecedent for our pinned-field and optimized-path
comparisons is Aronow and Lopatto~\cite{AL2026}. Hierarchical positive
laws, interpolation, and cavity calculations also have an established
lineage through Ruelle's cascade construction~\cite{Ruelle1987} and
the cavity variational principle of Aizenman, Sims, and
Starr~\cite{AizenmanSimsStarr2003}. Here we work with finite normalized
transition kernels and prove the quantitative larger-tree and
higher-moment estimates required by the size comparison.

These comparisons use normalized probability laws. An extra future
branch can be integrated out because its conditional law has mass one;
conditioning on an observation from that branch can instead change the
law retained in the calculation. PE's appendix ``Companion interfaces''
proves the polynomial-observation bound and permits an auxiliary common
temperature to vary with size; we recall these arguments when verifying
their use here. The changed-constraint estimates require the additional
local arguments developed in this article.
Appendix~\ref{sec:companion-interface} lists the imported estimates
and the hypotheses retained in their use.

Along a hierarchy we retain the initial disorder as a random input
and average over the later Gaussian increments. Ordinary averaging of
the resulting continuation gives one pressure; exponential averaging
with a small positive parameter $s$ gives the other. Their difference,
multiplied by $s$ times the number of sites, is the centered log Laplace
transform. We vary the covariance clocks and take finite differences
in the number of spins. Proposition~\ref{prop:increment} bounds their
combined effect with a power saving beyond the per-site scale $x^5$.
With $s=q x^{1+\varepsilon}$ and $N=x^{-6}$, multiplication by $Ns$
converts this into the saving in the displayed transform estimate.

The route from $N$ to $N'$ first changes the covariance data at fixed
size, then varies the size, and finally reverses the initial changes.
Writing $n=N e^\theta$ and $\ell=e^{\theta/6}$ keeps the endpoint
normalization equal to $n^{-1/6}$. Figure~\ref{fig:comparison-path}
shows these three stages.

\begin{figure}[htbp]
\centering
\begin{tikzpicture}[>=Latex,node distance=14mm,
 box/.style={draw=black!60,rounded corners=2pt,align=center,inner sep=7pt,text width=37mm,font=\small},
 lab/.style={font=\footnotesize,align=center,fill=white,inner sep=3pt}]
\node[box] (a) at (0,3.5) {Completed SK system at $N$\\$Y_N$};
\node[box] (b) at (8,3.5) {Completed SK system at $N'$\\$Y_{N'}$};
\node[box] (c) at (0,0) {Smoothed auxiliary\\system at $N$};
\node[box] (d) at (8,0) {Smoothed auxiliary\\system at $N'$};
\draw[->,thick] (a)--(c) node[midway,right,lab,text width=32mm] {Vary strength\\and fields;\\smooth the path};
\draw[->,thick] (c.south)--(0,-.9)--(8,-.9)--(d.south);
\node[lab,text width=105mm] at (4,-1.4) {Vary the clocks and compare neighboring sizes\\$\ell=e^{\theta/6}$};
\draw[->,thick] (d)--(b) node[midway,left,lab,text width=32mm] {Undo the smoothing,\\strength and field changes};
\draw[->,dashed] (a.north)--(0,4.5)--(8,4.5)--(b.north);
\node[lab,text width=105mm] at (4,4.9) {Resulting comparison of centered log Laplace transforms};
\end{tikzpicture}
\caption{The route from $N$ to an arbitrary $N'\in[N,2N]$. The displayed endpoints are the centered, rescaled bare systems. The proof bounds the changes along the route and the errors from the small regularizations at its ends. One parameter vector is selected for each finite route and each Laplace test.}
\label{fig:comparison-path}
\end{figure}


The comparison rests on four estimates.
\begin{enumerate}
\item We solve the linear equations for covariances of overlap errors.
The inverse bound remains uniform when the hierarchy is refined or its
masses become small. Summation by parts cancels the small-mass
denominators. Contributions from branching at an already specified
vertex are retained in the bound.
\item We vary the matrix and field covariances to turn pressure
comparisons into bounds on mean squared overlap errors. The constrained
minimization requires estimates in both the mass coordinate and the
quantile's value. A derivative of the minimum controls every minimizer,
so the argument does not need to select a minimizer measurably.
\item We obtain higher moments for the actual comparison paths by
projecting spin observations onto their revealed Gaussian histories and
integrating squared martingale increments. Before integrating a centered
quantity over a moving branch point, we show that the paths retained in
the calculation have a law independent of that point.
\item We sum the signed pressure coefficients before taking absolute
values. Summing the scalar contributions above a fixed cut, combining
the quadratic correction, and rescaling size cancel the leading terms
of degrees two, three, and four. The remaining increments are summable.
\end{enumerate}

\subsection{Consequences}

Theorem~\ref{thm:main} and the variance asymptotic above instantiate
the ferromagnetic comparison of Dey and Kang
\cite[Proposition~1.5]{DeyKang2026}. The comparison does not establish
the zero-field law, but it gives the following consequence once that
law and its positive variance prefactor are available.

\begin{corollary}[SK--Curie--Weiss fluctuations below unit coupling]
\label{cor:skcw}
Fix $\beta>1$ and $0\le\gamma<1$. Using the same independent standard
Gaussian couplings as in \eqref{eq:model}, define the exponent and log
partition function
\[
 \begin{gathered}
 \mathcal H_n^{\mathrm{CW}}(\sigma;\beta,\gamma)
 =
 \frac{\beta}{\sqrt n}\sum_{1\le i<j\le n}g_{ij}\sigma_i\sigma_j
 +\frac{\gamma}{2n}\left(\sum_{i=1}^n\sigma_i\right)^2,\\
 F_n^{\mathrm{CW}}(\beta,\gamma)
 =\log\sum_{\sigma\in\{-1,1\}^n}
 e^{\mathcal H_n^{\mathrm{CW}}(\sigma;\beta,\gamma)}.
 \end{gathered}
\]
Here $\gamma$ is the ferromagnetic coupling already in the exponent.
As $n\to\infty$ through all integers,
\[
 \begin{gathered}
 \Var F_n^{\mathrm{CW}}(\beta,\gamma)
 \sim a_2(\beta)n^{1/3},\\
 \frac{F_n^{\mathrm{CW}}(\beta,\gamma)
       -\mathbb E F_n^{\mathrm{CW}}(\beta,\gamma)}
      {\sqrt{\Var F_n^{\mathrm{CW}}(\beta,\gamma)}}
 \Longrightarrow\nu_\beta.
 \end{gathered}
\]
Here $a_2(\beta)>0$ is the constant $a_2$ above at the same $\beta$.
The centering and standard deviation are the perturbed model's own
exact finite-size quantities.
\end{corollary}

The comparison controls the added free energy in $L^2$ uniformly in
$n$. Since the unperturbed standard deviation grows like $n^{1/6}$,
this error disappears after either model's exact standardization.
The full argument is in Section~\ref{sec:consequences}.

This consequence is confined to the displayed zero-field ordered
Gaussian model at fixed $\beta>1$ and $0\le\gamma<1$; it does not cover
$\gamma=1$ or a coupling varying with $n$.

For the original zero-field model, the variance asymptotic also gives
the following consequence of Chatterjee's general Gaussian and
replacement inequalities.

\begin{corollary}[Upper disorder-chaos scales in zero-field SK]
\label{cor:disorder-chaos}
Fix $\beta>1$ and $n\ge2$. Put $E_n=\{(i,j):1\le i<j\le n\}$ and
$M=|E_n|=n(n-1)/2$, and let $g'$ be an independent copy of
$g=(g_{ij})_{(i,j)\in E_n}$ in \eqref{eq:model}. Write
$\langle\cdot\rangle_{g,h}$ for expectation under the product of that
model's Gibbs laws with disorders $g,h$, and
$R=n^{-1}\sum_i\sigma_i\tau_i$ for the overlap. Thus the replicas are
independent conditional on their disorders; $\mathbb E$ below also
averages the disorder and any replacement set.

For $t>0$, set $g^t=e^{-t}g+\sqrt{1-e^{-2t}}\,g'$. For an integer
$0\le k\le M$, let $A$ be independent of $g,g'$ and uniform over the
$k$-element subsets of $E_n$. Set $g^A_{ij}=g'_{ij}$ on $A$ and
$g^A_{ij}=g_{ij}$ otherwise. Then
\begin{align}
 \mathbb E\langle R^2\rangle_{g,g^t}
 &\le \frac1n+\frac{2\Var F_n(\beta)}
                 {\beta^2n(1-e^{-t})},\label{eq:chaos-ou}\\
 \mathbb E\langle R^2\rangle_{g,g^A}
 &\le \frac1n+\frac{2(M+1)\Var F_n(\beta)}
                 {\beta^2n(k+1)}+C_\beta n^{-1/2},\label{eq:chaos-replacement}
\end{align}
for a finite constant $C_\beta$ depending only on $\beta$. As $n\to\infty$, for
$t_n\ge0$ and integers $0\le k_n\le M$, these bounds imply
\[
 \begin{aligned}
 t_n n^{2/3}\to\infty
 &\quad\Longrightarrow\quad
 \mathbb E\langle R^2\rangle_{g,g^{t_n}}\to0,\\
 (k_n/M)n^{2/3}\to\infty
 &\quad\Longrightarrow\quad
 \mathbb E\langle R^2\rangle_{g,g^{A_n}}\to0,
 \end{aligned}
\]
where $A_n$ has size $k_n$ as above.
\end{corollary}

Chatterjee's variance identity and replacement inequality bound the
mean squared overlap by the free-energy variance. Substitution of the
$n^{1/3}$ variance asymptotic gives the displayed sufficient scales;
Section~\ref{sec:consequences} checks both normalizations.

These are sufficient annealed perturbation scales, not a sharp
transition or a quenched statement. No overlap conclusion is asserted
for the SK--Curie--Weiss model in Corollary~\ref{cor:skcw}.

\subsection{Organization and order of limits}

Sections~\ref{sec:foundations}--\ref{sec:paths} develop the finite-tree
calculus and admissible comparison paths. The pair inverse in
Section~\ref{sec:inverse} first supplies the global trial accuracy of
Section~\ref{sec:global}; that accuracy and the local auxiliary comparisons
give the genuine-path budgets of Section~\ref{sec:budgets}. The analytic
inputs for the modified optimizers are proved within that section before
their use. Section~\ref{sec:statistics} derives genuine-path statistics
from the budgets and uses the pair inverse again in its high-mass
covariance equation. Scalar compression in Section~\ref{sec:expansion}
and the common parameter order in Section~\ref{sec:parameters} complete
Proposition~\ref{prop:increment}. The conditional deduction in
Section~\ref{sec:reduction} then gives full-sequence moment convergence,
positive limiting variance, moment determinacy, the finite-system law
specification, and exact finite-size standardization. Section~\ref{sec:consequences}
contains the proofs of the two corollaries stated above.

All conclusions are for a fixed $\beta>1$. Moment orders, expansion orders, and propagation depths are finite and fixed before $N\to\infty$. Their values may depend on the desired estimate. The exponent used for the small Laplace test is chosen only after the power saving in the increment estimate is established.
